\documentclass[11pt]{article}
\usepackage[T1]{fontenc}
\usepackage{microtype}
\begin{document}
\section{Margins and expansion}
Character protrusion lets punctuation, hyphens and the rounded sides of letters hang a
little into the margin, so that the edge of the text looks straight to the eye; font
expansion stretches or shrinks the letters of a line by a few thousandths, so that the
spaces between words vary less from line to line. ``Quotes,'' hyphen-ated words, and
full stops. Both change where lines break, since a line can take a little more or a
little less text. \textit{Italic text, with commas, and a few more words that wrap} and
\textbf{bold text, which breaks onto a second line when it is long enough to do so}, and
\texttt{typewriter text, which protrudes but does not expand}. \small Small text at the
end, which also wraps to a further line to show the smaller size, with a final stop.
\end{document}
